\documentclass[10pt,a4paper]{amsart}
\usepackage[margin=19mm]{geometry}
\usepackage{amsmath,amssymb,mathtools}
\usepackage[scr=rsfs,cal=euler]{mathalfa}
\usepackage{xfrac}
\usepackage[unicode,hidelinks,bookmarks=false]{hyperref}
\newcommand{\ie}{i.e.\ }
\newcommand{\cf}{cf.\ }
\newcommand{\resp}{respectively\ }
\newcommand{\Subsection}{\subsection}
\providecommand{\nobreakdash}{\nobreak\hbox{-}}
\setlength{\emergencystretch}{2em}
\multlinegap0pt
\usepackage{amssymb}
\newtheorem{lemma}{Lemma}
\newtheorem{corollary}{Corollary}
\newcommand{\R}{\mathbb{R}}
\newcommand{\ip}[2]{\langle #1,#2\rangle}
\title{Counterexamples for {BFGS}-type methods under arbitrary strong {Wolfe} constants}
\author{Rui Diao}
\date{Source: arXiv:2609.09686v1 (CC BY 4.0)}
\begin{document}
\maketitle

\noindent\emph{Source and licence.} Counterexamples for {BFGS}-type methods under arbitrary strong {Wolfe} constants. Version arXiv:2609.09686v1 (CC BY 4.0), distributed by arXiv under the Creative Commons Attribution 4.0 International licence, http://creativecommons.org/licenses/by/4.0/, as recorded in the arXiv licence metadata for this exact version and shown on the abstract page https://arxiv.org/abs/2609.09686v1. Full licence notice: \texttt{licenses/arxiv-2609.09686-CC-BY-4.0.txt}.\par\smallskip

\noindent\textit{Editorial scope.} Counted unit: the article's lemma lem:global (source0.tex lines 993--1001). The article's complete original proof of it is delivered with it (source0.tex lines 1003--1102, one \texttt{proof} environment of 100 lines). No mathematics is written, repaired or re-derived: the statement and the proof are the article's own lines, in the article's own notation, and no retained line is substituted or re-flowed.  Retained uncounted as the source-local context the proof needs: lines 399--402; lines 431--456; lines 474--479; lines 520--523; lines 536--539; lines 566--570; lines 577--580; lines 589--592; lines 596--599; lines 692--705; lines 738--746; lines 764--773; lines 847--858; lines 878--882; lines 953--956; lines 960--964; lines 985--991; lines 1293--1302.  Nothing was omitted from any retained span, so the package carries no omission record.  No retained line contains a citation, so the wrapper carries no bibliography and no bibliography entry.  No retained line contains a figure environment, an image-inclusion command or drawing code, so no image is delivered and the package declares no asset dependency.  Editorial formatting only, in addition: an AMS \texttt{amsart} preamble that restates the article's own declarations and macros used by the retained lines, namely the environments \texttt{lemma}, \texttt{corollary} on separate counters, together with the standard packages those lines need.  The retained environments print in this document as Lemma 1, Corollary 1 in order of appearance, whereas the article prints the same items with the numbering of its own full document.  The complete native source of the article as distributed in the arXiv e-print for this exact version is bundled unchanged as \texttt{sources/209845/source0.tex}.
\begin{equation}
\label{eq:S0}
\mathcal S_0 = \{x_0 + \tau s_0 + n\nu_0 : \tau \in (-\delta, 1],\, |n| < \delta\}.
\end{equation}

\begin{lemma}[Overlaps of strips and disks]
\label{lem:seam}
Let $0 < \delta \le 1/20$.
\begin{enumerate}
\item[(i)] \emph{Non-adjacent strips and disks are separated.}
Non-adjacent strips and disks have disjoint closures:
\[
\overline{\mathcal S_i} \cap \overline{\mathcal S_j} = \emptyset \ \ (|i-j| \ge 2), \qquad
\overline{\mathcal B_k} \cap \overline{\mathcal S_j} = \emptyset \ \ (j \notin \{k-1, k\}), \qquad
\overline{\mathcal B_k} \cap \overline{\mathcal B_j} = \emptyset \ \ (j \ne k).
\]
Moreover every point of $\R^2$ has a neighborhood meeting at most two strips (necessarily consecutive) and at most one vertex disk.
\item[(ii)] \emph{Overlaps at a corner lie in the collar.}
For every $k \ge 1$, all overlaps among the adjacent strips $\mathcal S_{k-1}, \mathcal S_k$ and the vertex disk $\mathcal B_k$ lie in the open Fermi wedge, and hence in the collar $\mathcal C_k$:
\begin{equation}
\label{eq:cornerincl}
(\mathcal S_{k-1} \cap \mathcal S_k) \cup \bigl(\mathcal B_k \cap (\mathcal S_{k-1} \cup \mathcal S_k)\bigr) \subset \{\tau_{k-1}>1-\delta\}\cap\{\tau_k<\delta\}.
\end{equation}
Moreover, restricting $\mathcal C_k$ to either strip recovers precisely that strip's own collar:
\begin{equation}
\label{eq:Ckrest}
\mathcal C_k \cap \mathcal S_{k-1} = \{p \in \mathcal S_{k-1} : \tau_{k-1}(p) \ge 1-\delta\}, \qquad
\mathcal C_k \cap \mathcal S_k = \{p \in \mathcal S_k : \tau_k(p) \le \delta\}.
\end{equation}
\end{enumerate}
\end{lemma}

\begin{equation}
\label{eq:corner}
\operatorname{dist}(p, \ell_{k-1}) < \delta \implies \tau_k(p) < \delta,
\qquad
\operatorname{dist}(p, \ell_k) < \delta \implies \tau_{k-1}(p) > 1-\delta.
\end{equation}

\begin{equation}
\label{eq:Qk}
Q_k(x) = f_k + \ip{g_k}{x-x_k} + \frac{\alpha_k}{2}\|x-x_k\|^2.
\end{equation}

\begin{equation}
\label{eq:template}
F_k(\tau, n) = a(\tau) + b(\tau)n + \tfrac12 a''(\tau)n^2.
\end{equation}

\begin{equation}
\label{eq:btrans}
b(\tau) = \ip{g_k}{\nu_k} + \bigl(\ip{g_{k+1}}{\nu_k} - \ip{g_k}{\nu_k}\bigr)
S(\tau; \delta, 1-\delta).
\end{equation}

\begin{equation}
\label{eq:Adef}
a''(\tau) = \Lambda_k^-(\tau) + \Lambda_k^+(1-\tau),
\end{equation}

\begin{equation}
\label{eq:adef}
a_k(0)=f_k, \qquad a_k'(0)=\eta_k, \qquad f_{k+1}:=a_k(1) \qquad (k\ge0);
\end{equation}

\begin{equation}
\label{eq:Q0collar}
F_0(\tau, n) = -\tau = Q_0(\tau, n) \quad\text{on } (-\delta, \delta] \times (-\delta, \delta).
\end{equation}

\begin{lemma}[Axial profile]
\label{lem:profile}
Fix $k\ge0$ and $\lambda\in[0,1]$, and let $a''$ denote the axial curvature \eqref{eq:Adef} with this split parameter. Then:
\begin{enumerate}
\item[(i)] \emph{Strict axial descent and stationarity.} $a'(1)=0$ and $a'<0$ on $[0,1)$; consequently $a$ is strictly decreasing on $[0,1]$, with $f_{k+1}<f_k$.
\item[(ii)] \emph{Collar quadratic forms.} $a(\tau)=f_k-\alpha_{k+1}\tau+\tfrac12\alpha_k\tau^2$ on $[0,\delta]$, and $a(\tau)=f_{k+1}+\tfrac12\alpha_{k+1}(\tau-1)^2$ on $[1-\delta,1]$.
\item[(iii)] \emph{Net descent and the Armijo ratio.}
\begin{equation}
\label{eq:rhocentroid}
f_k-f_{k+1}=\int_0^1\tau\,a''(\tau)\,d\tau, \qquad
\frac{f_{k+1}-f_k}{\eta_k} = \frac1{\alpha_{k+1}}\int_0^1\tau\,a''(\tau)\,d\tau.
\end{equation}
\end{enumerate}
\end{lemma}

\begin{lemma}[Collar agreement and smoothness]
\label{lem:collar}
For any $k\ge0$ and any $\lambda_k\in[0,1]$, the formulas \eqref{eq:btrans} and \eqref{eq:Adef} define $b,a''\in C^\infty(\R)$, and with $a,a'$ obtained by integrating from the left-vertex data \eqref{eq:adef}, the function
\begin{equation}
\label{eq:Ftilde}
\widetilde F_k(x) := a(\tau_k)+b(\tau_k)n_k+\tfrac12 a''(\tau_k)n_k^2
\end{equation}
belongs to $C^\infty(\R^2)$. Moreover $\widetilde F_k\equiv Q_k$ on the half-plane $\{\tau_k\le\delta\}$ and $\widetilde F_k\equiv Q_{k+1}$ on the half-plane $\{\tau_k\ge 1-\delta\}$. In particular, the restriction $F_k=\widetilde F_k\big|_{\mathcal S_k}$ is $C^\infty$ on $\mathcal S_k$ and coincides with $Q_k$ on the left collar and with $Q_{k+1}$ on the right collar. For $k=0$, $Q_0$ is linear and the left identity holds on $\{\tau_0\le\delta\}$, hence throughout the backward tail of \eqref{eq:S0}.
\end{lemma}

\begin{lemma}[Uniform Hessian bound]
\label{lem:hessU}
For any $k\ge0$ and any $\lambda_k\in[0,1]$, the Hessian of the strip template $F_k$ in \eqref{eq:template} satisfies
\begin{equation}
\label{eq:hessU}
\|\nabla^2 F_k(\tau,n)\|\le\frac{K_{\mathcal U}}{\delta}
\quad\text{for all }(\tau,n)\in\mathcal S_k,
\end{equation}
where $K_{\mathcal U} > 0$ is a universal constant independent of $k$ and $\lambda_k$.
\end{lemma}

\begin{corollary}[Well-defined interpolant on $\mathcal U$]
\label{cor:F}
The function $F: \mathcal U \to \R$, defined on each piece of $\mathcal U = \bigcup_{k\ge 1}\mathcal B_k \cup \bigcup_{k\ge 0}\mathcal S_k$ by
\begin{equation}
\label{eq:Fdef}
F(x) = \begin{cases}
F_k(\tau_k, n_k), & x \in \mathcal S_k \quad (k \ge 0), \\
Q_k(x), & x \in \mathcal B_k \quad (k \ge 1),
\end{cases}
\end{equation}
is well-defined and belongs to $C^\infty(\mathcal U)$.
\end{corollary}

\begin{lemma}[Exterior limit value]
\label{lem:finf}
$(f_k)_{k\ge0}$ is strictly decreasing, and $f_\infty:=\inf_kf_k$ is
finite, with $f_k\downarrow f_\infty$.
\end{lemma}

\begin{equation}
\label{eq:Phidef}
\Phi(x) = \sum_{k\ge 1} V_k(x) + \sum_{k\ge 0} A_k(\tau_k) W(n_k), \qquad \zeta(x) = S(\Phi(x)).
\end{equation}

\begin{lemma}[Core plateau and support]
\label{lem:core}
On the polyline $P$ and on each inner vertex disk $B(x_k, r_\circ)$ ($k \ge 1$), $\Phi \ge 1$, $\zeta \equiv 1$, and $\nabla\zeta \equiv 0$.
Moreover, $\operatorname{supp}\zeta \subset \overline{\mathcal U}$.
\end{lemma}

\begin{equation}
\label{eq:fdef}
f(x) = \begin{cases}
\zeta(x) F(x) + (1-\zeta(x))f_\infty, & x \in \mathcal U, \\
f_\infty, & x \in \R^2\setminus \mathcal U.
\end{cases}
\end{equation}

\begin{lemma}[Smooth switch $S$]
\label{lem:S}
$S\in C^\infty(\R)$, and:
\begin{enumerate}
\item[(i)] $S(u)=0$ for $u\le0$, $S(u)=1$ for $u\ge1$, and $S(u)\in[0,1]$ for every $u\in\R$;
\item[(ii)] $S^{(m)}(u)=0$ for every $u\notin(0,1)$ and every $m\ge1$;
\item[(iii)] $\int_0^1S(u)\,du=\tfrac12$;
\item[(iv)] every derivative of $S$ is bounded on $\R$.
\end{enumerate}
\end{lemma}

\begin{lemma}[Global extension and regularity]
\label{lem:global}
The function $f$ in \eqref{eq:fdef} satisfies:
\begin{enumerate}
\item[(i)] $f\in C^\infty(\R^2)$, and along the polyline $P$, $f\equiv F$ and $\nabla f=\nabla F$;
\item[(ii)] $f$ is bounded below: $f(x)\ge f_\infty-2\delta$ for all $x\in\R^2$;
\item[(iii)] $\nabla f$ is Lipschitz continuous on $\R^2$, with $\|\nabla^2f\|\le K/\delta^2$ for a universal constant $K > 0$.
\end{enumerate}
\end{lemma}

\begin{proof}
(i) \emph{Smoothness and polyline agreement.}
The function $\zeta$ belongs to $C^\infty(\R^2)$, and $\zeta\equiv 0$ on $\R^2\setminus\overline{\mathcal U}$ by Lemma~\ref{lem:core}, so $f\equiv f_\infty$ is smooth on the open exterior.
On $\mathcal U$, $F$ is $C^\infty$ (Corollary~\ref{cor:F}) and $\zeta$ is $C^\infty$, so $f$ is smooth on $\mathcal U$.
It remains to check the interface $\partial\mathcal U$.

Let $p\in\partial\mathcal U$. By Lemma~\ref{lem:seam}(i), some neighborhood of $p$ meets at most two consecutive strips and at most one vertex disk. Shrinking that neighborhood to a ball $V$ about $p$, we produce $G\in C^\infty(V)$ with $F=G$ on $\mathcal U\cap V$, whence $f=f_\infty+\zeta(G-f_\infty)$ on $V$.

If $p$ lies in $\overline{\mathcal S_{k-1}}\cap\overline{\mathcal S_k}$ for some $k\ge 1$, then $\operatorname{dist}(p,\ell_{k-1})\le\delta$ and $\operatorname{dist}(p,\ell_k)\le\delta$, so $\tau_k(p)\le\delta$ and $\tau_{k-1}(p)\ge 1-\delta$ by the same comparison as \eqref{eq:corner}, now with non-strict inequalities.
With $q$ and $b$ as in the proof of \eqref{eq:corner}, the comparison $\tau_k(p)=\ip{p-q}{s_k}-b\ip{s_{k-1}}{s_k}\le\operatorname{dist}(p,\ell_{k-1})$ becomes an equality only if $b=0$ and $p-q$ is a nonnegative multiple of $s_k$, hence only if $p=x_k+\delta s_k$, which lies in the interior of $\mathcal U$; symmetrically, $\tau_{k-1}(p)=1-\delta$ would force $p=x_k-\delta s_{k-1}$, likewise interior.
Thus $p$ lies in the open wedge $\{\tau_{k-1}>1-\delta\}\cap\{\tau_k<\delta\}$. Shrinking $V$ into this wedge, Lemma~\ref{lem:collar} gives $\widetilde F_{k-1}\equiv Q_k\equiv\widetilde F_k$ on $V$, so $F=Q_k$ on $\mathcal U\cap V$. Take $G=Q_k$.

If $p$ lies in exactly one closed strip $\overline{\mathcal S_j}$, shrink $V$ to miss every other strip (Lemma~\ref{lem:seam}(i)). Any disk met by $V$ is incident to $\mathcal S_j$: $x\in\mathcal B_j$ implies $\tau_j(x)\le\|x-x_j\|<\delta$, while $x\in\mathcal B_{j+1}$ implies $1-\tau_j(x)\le\|x-x_{j+1}\|<\delta$, so $\widetilde F_j\equiv Q_j$ on $\mathcal B_j$ and $\widetilde F_j\equiv Q_{j+1}$ on $\mathcal B_{j+1}$ by Lemma~\ref{lem:collar}. Thus $F=\widetilde F_j$ on $\mathcal U\cap V$. Take $G=\widetilde F_j$.

If $p$ lies in no closed strip, then $p\in\partial\mathcal B_k$ for some $k\ge 1$. Shrinking $V$ so that $\mathcal U\cap V\subset\mathcal B_k$, one has $F=Q_k$ on $\mathcal U\cap V$. Take $G=Q_k$.

In all cases $f$ is $C^\infty$ near $p$, so $f\in C^\infty(\R^2)$.

Along the polyline $P \subset \mathcal U$, Lemma~\ref{lem:core} gives $\zeta \equiv 1$ and $\nabla\zeta \equiv 0$.
Differentiating $f = f_\infty + \zeta(F - f_\infty)$ along $P$ gives
\[
f = F, \qquad \nabla f = \zeta\nabla F + (F - f_\infty)\nabla\zeta = \nabla F.
\]

(ii) \emph{Lower bound.}
On $\mathcal U$, $f = \zeta F + (1-\zeta)f_\infty$ is a convex combination of $F$ and $f_\infty$ (since $\zeta \in [0, 1]$).
Because $f_\infty \ge f_\infty - 2\delta$ trivially, it suffices to show $F \ge f_\infty - 2\delta$ on each piece of $\mathcal U$:
\begin{itemize}
\item \emph{Vertex disks $\mathcal B_k$ ($k\ge 1$):}
Dropping the non-negative quadratic term $\frac12\alpha_k\|x-x_k\|^2$ in $Q_k$ \eqref{eq:Qk} gives
\[
Q_k(x) \ge f_k + \ip{g_k}{x-x_k} \ge f_k - \|g_k\|\|x-x_k\| = f_k - \|x-x_k\|,
\]
since $\|g_k\| = 1$.
For $x \in \mathcal B_k = B(x_k, \delta)$, $\|x-x_k\| < \delta$.
Since $f_k \ge f_\infty$ (Lemma~\ref{lem:finf}),
\[
Q_k(x) > f_k - \delta \ge f_\infty - \delta > f_\infty - 2\delta.
\]

\item \emph{Segment strips $\mathcal S_k$ ($k\ge 0$):}
For every $k\ge 0$, the function on $\mathcal S_k$ is given by the strip template $F_k(\tau, n) = a(\tau) + b(\tau)n + \frac12 a''(\tau)n^2$.
By Lemma~\ref{lem:profile}(i), $a$ is strictly decreasing on $[0, 1]$ (for $k=0$, this decrease extends back onto $[-\delta, 0]$ where $a' \equiv -1$).
Hence throughout $\mathcal S_k$, $a(\tau) \ge a(1) = f_{k+1} \ge f_\infty$.
The transverse tilt $b(\tau)$ is a convex combination of $\ip{g_k}{\nu_k}$ and $\ip{g_{k+1}}{\nu_k}$, so $|b(\tau)| \le 1$, giving $b(\tau)n \ge -\delta$ since $|n| < \delta$.
Because the axial curvature satisfies $a'' \ge 0$, the quadratic term $\frac12 a''(\tau)n^2 \ge 0$ is non-negative.
Combining these estimates yields
\[
F_k(\tau, n) \ge f_{k+1} - \delta \ge f_\infty - \delta > f_\infty - 2\delta.
\]
\end{itemize}
Therefore, $F \ge f_\infty - 2\delta$ everywhere on $\mathcal U$.
Taking the convex combination, $f(x) \ge f_\infty - 2\delta$ for all $x \in \R^2$.

(iii) \emph{Lipschitz gradient bound.}
Differentiating the global interpolant $f = f_\infty + \zeta(F - f_\infty)$ gives the gradient vector field $\nabla f = \zeta \nabla F + (F - f_\infty)\nabla\zeta$.
Differentiating once more, the spatial Hessian on $\mathcal U$ is given pointwise by
\begin{equation}
\label{eq:hess-pointwise}
\nabla^2 f = \zeta \nabla^2 F + \nabla\zeta \otimes \nabla F + \nabla F \otimes \nabla\zeta + (F - f_\infty)\nabla^2\zeta.
\end{equation}
Outside $\overline{\mathcal U}$, $\zeta \equiv 0$ identically, so $\nabla^2 f \equiv 0$.
Taking spectral norms on $\mathcal U$, the triangle inequality gives
\begin{equation}
\label{eq:hess-norm-bound}
\|\nabla^2 f\| \le \|\zeta\|_\infty \|\nabla^2 F\| + 2\|\nabla\zeta\|\|\nabla F\| + \|F - f_\infty\|_\infty \|\nabla^2\zeta\|.
\end{equation}
We now bound each factor across $\mathcal U$:
\begin{enumerate}
\item \emph{Regularity and amplitude of $F$:}
By Lemma~\ref{lem:hessU} on each strip and $\nabla^2 Q_k = \alpha_k I$ ($\alpha_k \le 1$) on vertex disks, $\sup_{\mathcal U}\|\nabla^2 F\| \le K_{\mathcal U}/\delta = O(1/\delta)$.
On the backward tail $[-\delta, 0] \times (-\delta, \delta)$ around $x_0$, $F_0(\tau, n) = -\tau \le \delta$.
Because $(f_k)$ is strictly decreasing with $f_0 = 0$ (Lemma~\ref{lem:finf}), $a_k(\tau) \le f_k \le f_0 = 0$ for $\tau \in [0, 1]$ on each strip $\mathcal S_k$.
Combined with the transverse bound $|b(\tau)n| + \frac12 a''(\tau)n^2 \le \delta + \frac12\delta < 2\delta$ (since $a'' < 1/\delta$ by Lemma~\ref{lem:hessU}), $F_k \le a_k(\tau) + 2\delta \le 2\delta$, while on vertex disks $Q_k \le f_k + \delta + \frac12\delta^2 < 2\delta$.
Together with the lower bound $F \ge f_\infty - 2\delta$ from part~(ii), $F(x) \in [f_\infty - 2\delta, 2\delta]$ across $\mathcal U$, and consequently $\|F - f_\infty\|_{L^\infty(\mathcal U)} \le |f_\infty| + 2\delta = O(1)$.
On the backward extension $(-\delta, 0] \times (-\delta, \delta)$, $F_0(\tau, n) = -\tau$ by \eqref{eq:Q0collar} is linear with $\|\nabla F_0\| = \|s_0\| = 1$.
On the remainder of $\mathcal U$, each point $x$ connects to the polyline $P$ via a transverse or radial segment of length at most $2\delta$ within $\mathcal U$.
Along $P$, $\|\nabla F\| \le \sqrt2$, so the mean value theorem gives
\[
\|\nabla F\|_\infty \le \max\bigl(1, \sqrt2 + \|\nabla^2 F\|(2\delta)\bigr) \le \sqrt2 + 2K_{\mathcal U} = O(1).
\]

\item \emph{Cutoff derivatives:}
Each constituent bump in \eqref{eq:Phidef} is defined via the smooth transition $S$ on intervals of length scale $\delta$.
Summing at most three overlapping bumps gives $\|\nabla\Phi\|_\infty = O(1/\delta)$ and $\|\nabla^2\Phi\|_\infty = O(1/\delta^2)$.
With $\zeta = S(\Phi)$ and $\|S'\|_\infty, \|S''\|_\infty < \infty$ (Lemma~\ref{lem:S}(iv)), the chain rule gives $\|\nabla\zeta\|_\infty \le \|S'\|_\infty \|\nabla\Phi\|_\infty = O(1/\delta)$, and
\[
\nabla^2\zeta = S''(\Phi)\nabla\Phi\otimes\nabla\Phi + S'(\Phi)\nabla^2\Phi,
\]
which yields $\|\nabla^2\zeta\|_\infty \le \|S''\|_\infty \|\nabla\Phi\|_\infty^2 + \|S'\|_\infty \|\nabla^2\Phi\|_\infty \le O(1/\delta^2) + O(1/\delta^2) = O(1/\delta^2)$.
\end{enumerate}
Substituting these bounds into \eqref{eq:hess-norm-bound}, and using $\|\zeta\|_\infty \le 1$ (Lemma~\ref{lem:S}(i)), gives
\[
\|\nabla^2 f(x)\| \le O(1/\delta) + O(1/\delta) + O(1/\delta^2) \le \frac{K}{\delta^2} \quad \text{for all } x \in \R^2,
\]
where $K > 0$ is a universal constant.
Because $\R^2$ is convex and $f \in C^\infty(\R^2)$, the mean value theorem yields
$\|\nabla f(x) - \nabla f(y)\| \le (K/\delta^2)\|x-y\|$ for all $x, y \in \R^2$,
so $\nabla f$ is Lipschitz continuous on $\R^2$.
\end{proof}

\end{document}
